
\providecommand{\appendixDir}{Sections/appendix}

\makeatletter
\@ifundefined{clawcode}{%
  \GenericError{}{%
    UniClawBench appendix: the `clawcode' environment is undefined.\MessageBreak
    Add the following line to your main paper preamble (above\MessageBreak
    \string\begin{document}):\MessageBreak
    \space\space\string\usepackage{tcolorbox}
\usepackage{makecell}
\tcbuselibrary{listings, breakable, skins}

\newtcblisting{clawcode}[1][]{%
  listing only,
  breakable,
  enhanced,
  colback=black!2,
  colframe=black!55,
  boxrule=0.4pt,
  arc=2pt,
  outer arc=2pt,
  left=6pt, right=6pt, top=4pt, bottom=4pt,
  fonttitle=\bfseries\small,
  coltitle=black,
  attach boxed title to top left={xshift=4pt, yshift=-3pt},
  boxed title style={%
    colback=white,
    colframe=black!55,
    sharp corners=northwest,
    arc=0pt,
    boxrule=0.3pt,
    top=1pt, bottom=1pt, left=4pt, right=4pt,
  },
  listing options={%
    basicstyle=\ttfamily\small,
    breaklines=true,
    breakatwhitespace=true,
    breakautoindent=false,
    breakindent=0pt,
    prebreak={},
    postbreak={},
    columns=fullflexible,
    keepspaces=true,
    upquote=true,
    showstringspaces=false,
    aboveskip=4pt, belowskip=4pt,
  },
  #1
}

\newtcblisting{clawsnippet}{%
  listing only,
  breakable,
  enhanced,
  colback=black!2,
  colframe=black!40,
  boxrule=0.3pt,
  arc=1.5pt,
  left=5pt, right=5pt, top=3pt, bottom=3pt,
  listing options={%
    basicstyle=\ttfamily\footnotesize,
    breaklines=true,
    columns=fullflexible,
    keepspaces=true,
    upquote=true,
    showstringspaces=false,
  },
}

\newtcolorbox{clawschema}[1][]{%
  breakable,
  enhanced,
  colback=black!2,
  colframe=black!55,
  boxrule=0.4pt,
  arc=2pt,
  outer arc=2pt,
  left=6pt, right=6pt, top=4pt, bottom=4pt,
  fonttitle=\bfseries\small,
  coltitle=black,
  attach boxed title to top left={xshift=4pt, yshift=-3pt},
  boxed title style={%
    colback=white,
    colframe=black!55,
    sharp corners=northwest,
    arc=0pt,
    boxrule=0.3pt,
    top=1pt, bottom=1pt, left=4pt, right=4pt,
  },
  #1
}

\newcommand{\schemaitem}[2]{%
  \par\smallskip
  \noindent\texttt{\bfseries\#\# #1}\par
  \nopagebreak
  \noindent{\itshape\color{black!60}\small #2}\par
}
%
  }{%
    The preamble file loads tcolorbox and defines the boxed-listing\MessageBreak
    environments used by sections A--D. See appendix/README.md.%
  }{}%
  \endinput
}{}
\makeatother


\section{Task Definition Format}\label{app:format}

A UniClawBench task is generally defined by two files: a single task
declaration written in YAML that fixes the public contract, and a
hidden \texttt{eval\_rule.md} that specifies the rubric. The executor
sees only the \texttt{task} field rendered into a runtime prompt; the
supervisor sees both files and treats the eval rule as the canonical
scoring contract. Appendix~\ref{app:format:yaml} gives the YAML template that fixes
the public contract; Appendix~\ref{app:format:fields} walks through the
semantics of each field; and Appendix~\ref{app:format:eval} specifies the
nine-section schema of the hidden eval rule.

\subsection{Task YAML Template}\label{app:format:yaml}

\begin{clawcode}[title={Task YAML Template},label={lst:format:yaml}]
task_id: task_NNN_short_slug
category: <suite_name>
agent_sys: openclaw
agent_id: main
model: claude-opus-4.6

timeout_seconds: 1200
max_total_seconds: 1800
success_threshold: 0.90

task: |
  <natural-language request shown to the executor>

task_snapshot: |
  <parallel-safe variant>

references:
  - references/eval_rule.md
  - references/ground_truth.json
  - ...

sources:
  - <relative_path>

skills:
  - <skill_id>

services:
  - name: <service>
    path: <subdir>
    start: bash install.sh
    oneshot: true

pre_exec:
  - ops/populate.py
pre_exec_parallel_safe: true

codex:
  max_user_followups: 2
  user_simulator:
    policy: |
      <override of the default user-simulator policy>
  supervisor:
    instructions: |
      <task-specific addendum to the default supervisor instructions>
\end{clawcode}

\subsection{Field Semantics}\label{app:format:fields}

\paragraph{Basic definition.}
\texttt{task\_id} is unique within the suite and is also the file
stem. \texttt{category} is one of the five task dimensions from the main
paper. \texttt{agent\_sys}, \texttt{agent\_id}, and \texttt{model}
record the \emph{default} executor binding when the task is
dispatched without an override; cross-architecture and cross-model
sweeps override these at runtime.

\paragraph{Lifecycle.}
\texttt{timeout\_seconds} bounds a single executor cycle.
\texttt{max\_total\_seconds} bounds the whole-attempt wall clock
across all cycles, excluding supervisor and user-simulator turns.
\texttt{success\_threshold} is the minimum supervisor score that
records the attempt as \texttt{Pass}. The cycle count is bounded separately by
\texttt{codex.max\_user\_followups}.

\paragraph{Task prompt.}
\texttt{task} is the literal natural-language request the executor
receives, wrapped in a small runtime preamble that surfaces workspace
paths and installed skills. \texttt{task\_snapshot} is an optional
snapshot version that consumes snapshot data
instead of a live API.

\paragraph{Injections.}
\texttt{sources} lists files placed under the executor's workspace mount.
\texttt{skills} lists declared skills; the eval rule for a
skill-usage task typically requires transcript evidence that the
declared skill was actually consulted. \texttt{services} entries are
docker-compose-style bootstrap services started before the executor's
first turn (databases, mock APIs, GUI app installers).
\texttt{pre\_exec} scripts run once after services are healthy and
typically populate session-specific fixture state into live API.
\texttt{references} are hidden judging assets placed under the
supervisor's workspace and never copied into the executor container.

\paragraph{Supervisor and user simulator.}
The \texttt{codex} block configures the closed-loop interaction.
\texttt{max\_user\_followups} bounds the number of follow-up cycles
allowed before the runner forcibly closes the attempt.
\texttt{codex.user\_simulator.policy} overrides the default
user-simulator behavior policy when a task needs a non-standard
interaction style (e.g.\ a deliberately terse user, a user with a
specific persona constraint).
\texttt{codex.supervisor.instructions} appends a task-specific
addendum to the default supervisor template, used to highlight
checkpoints that are easy to miss without explicit attention or to
narrow down what counts as evidence for an unusual artifact.
Both blocks are optional; when absent, the defaults from
Appendix~\ref{app:prompts} apply.

\subsection{Eval Rule (\texttt{eval\_rule.md}) Format}\label{app:format:eval}

Every task ships a hidden Markdown rubric at
\texttt{references/eval\_rule.md}. The rubric follows a fixed
nine-section schema: \S1--\S4 establish the contract, \S5--\S6
define scoring, and \S7--\S9 cover lifecycle policy and asset
bookkeeping. The supervisor prompt instructs the supervisor to score
using the lines that appear in \S5 and \S6, and other parts guide the
understanding and flexible judgment.

\begin{clawschema}[title={Eval rule schema},label={lst:format:eval}]
\schemaitem{1. Grading Philosophy}{States what the task tests, what
  is rewarded, and what should not be penalized. Notes the role of
  declared skills.}

\schemaitem{2. Task Contract}{Restates the public prompt with
  deliverable paths pinned, and separates hard requirements (must
  produce) from soft requirements (must follow these formats).}

\schemaitem{3. Source-Selection and Target-Resolution Rules}{How the
  supervisor maps executor artifacts to ground-truth entities --
  alias rules, filename-stem variants, numeric tolerance bands.}

\schemaitem{4. Ground-Truth Snapshot}{Pointer to ground-truth files
  with canonical entity counts and any must-not-appear entities for
  negative checkpoints.}

\schemaitem{5. Checkpoint Rubric}{Numbered scoring lines whose
  weights sum to 1.00. Each is boolean, graded, or count-based. The
  supervisor scores using only these lines.}

\schemaitem{6. Scoring Policy \& Score Caps}{Hard ceilings applied
  after the \S5 sum. A cap is a maximum and cannot raise a \S5 sum
  that is already lower.}

\schemaitem{7. Continue vs Fail Guidance}{Score banding for the
  \texttt{pass} / \texttt{continue} / \texttt{fail} verdicts. The
  mid-band lets the supervisor request one focused follow-up.}

\schemaitem{8. Hidden Reference Assets}{Inventory of
  supervisor-only files; reaffirms that none are surfaced to the
  executor or user simulator.}

\schemaitem{9. Dynamic Content Note}{Whether the task is offline
  (static fixture) or live (consumes an external page that may
  shift between captures), with reconciliation rules for live
  tasks.}
\end{clawschema}

\section{Task Examples}\label{app:examples}

This section lists representative task examples from each suite. For
every task dimension we show two prompts. Each prompt is the literal
natural-language \texttt{task} field that the executor sees at
runtime; the surrounding YAML schema follows
Appendix~\ref{app:format}.

\subsection{Skill Usage}\label{app:examples:skill}

Tasks here ship several declared skills and require the executor to actually
consult it. The agent's job is to discover the skill, follow its
\texttt{SKILL.md} entry-point, and apply the right tool to the
provided data.

\begin{clawcode}[title={\texttt{task\_05\_org\_chart\_mermaid}}]
Our CEO asked for a governance visualization she can drop into her
board deck, and I need to pull it from the proxy statement. The
excerpt lives at
/tmp_workspace/clawbench/sources/alphabet_def14a_excerpt.md -- use the
Mermaid diagramming skill to turn it into a clean org chart. She wants
to see the actual reporting hierarchy -- the solid lines from each
current director to the board and from each non-CEO executive officer
to the executive roster -- who sits on which committee across the
audit, compensation, governance, and executive committees, and the
dotted advisory or observer relationships that aren't formal seats.
We also need who chairs what, and any seats that are vacant, interim,
advisor-only, or recently filled.

Please render it as Mermaid using graph TD or flowchart TD in
/tmp_workspace/results/org.mmd, and produce an SVG rendering at
/tmp_workspace/results/org.svg. Use solid arrows for board and
executive roster relationships that are explicitly supported by the
excerpt, dotted arrows for committee memberships and any non-voting
advisor or observer relationships, and group each committee as a
subgraph. Distinguish committee chairs visually from regular members
-- shape or styling is fine. Treat each person as a single node; if
someone holds multiple roles, draw their relationships from that one
node rather than duplicating them. Anyone listed as recently-departed
should not appear as a current node.

Also save /tmp_workspace/results/governance_matrix.csv so my chief of
staff can review the extraction without reading Mermaid. Use columns:
person, source_role, governance_group_or_board, committee,
committee_role, edge_type, source_quote.
\end{clawcode}

\begin{clawcode}[title={\texttt{task\_19\_github\_activity\_digest}}]
I missed the last two weeks of activity on clawbench-activity-fixture
and need to catch up before the maintainer sync. Build me a 16-day
activity digest covering 2026-04-10 through 2026-04-25 inclusive --
use the GitHub skill in our workspace for the API calls. Authenticate
with the \$GITHUB_TOKEN environment variable and hit the live GitHub
REST API directly via the gh CLI or curl against api.github.com. Do
not rely on any local cache or pre-exported snapshot file under
/tmp_workspace/clawbench for this run; everything must come from a
fresh API pull covering merged PRs, closed-unmerged PRs, closed
issues, and commits in that window.

A few entries in this fixture repo may have been imported by an
automation account. When a PR or issue body, or a commit message,
includes an explicit pseudo-user or planned-date note, use that note
as the contributor and date for the digest instead of the automation
account's login or the rebuild timestamp.

Save /tmp_workspace/results/digest.md grouped by contributor. For
each contributor show: PRs merged, PRs closed without merge, issues
closed, and commit count. Add a Follow-up Flags section for activity
that looks risky or needs a maintainer follow-up, including every
closed-unmerged PR and any revert PR or revert commit. End with a
Highlights section listing the 5 most recent merged PRs or closed
issues in the window, newest first. When two highlight events share
the same date, list the higher-numbered activity first so number 27
comes before number 13; if the numbers also tie, list the issue
before the pull request.

Also save /tmp_workspace/results/activity_index.csv with columns
activity_type,number_or_sha,title_or_message,contributor,activity_date,status
covering every merged PR, closed-unmerged PR, closed issue, and
commit in the 16-day window so I can sort and filter it later. Cover
the full set of contributors and activity, not a partial sample.
\end{clawcode}

\subsection{Exploration}\label{app:examples:exploration}

Tasks here are open-ended audits or refactors over noisy and
potentially misleading inputs. The agent must investigate multiple
choices, reject incorrect alternatives, and document negative
evidence to distinguish real exploration from a guessed conclusion.

\begin{clawcode}[title={\texttt{task\_01\_zshrc\_prompt\_refactor}}]
Your goal is to refactor the injected legacy `.zshrc` into a
maintainable, low-noise, verifiable zsh configuration. The input file
is at `/tmp_workspace/sources/existing_zshrc`. This is not about
writing a new prompt snippet from scratch; you must preserve the real
behaviors from the original config and eliminate the conflicts it
already contains.

Start by auditing the legacy config: identify duplicate PATH entries,
duplicate PROMPT/RPROMPT definitions, `precmd` overrides, Oh My
Zsh/plugin load order issues, conda/nvm initialization placement, and
any fragments that may produce side effects in interactive vs.
non-interactive shells. Then produce the final `.zshrc`: preserve
aliases, plugins, PATH prepends, the conda hook, and nvm sourcing;
the prompt should be concise, showing the current path and git
status, but must not display username, hostname, time,
Python/Node/Ruby runtime, full deep paths, or Nerd Font-only glyphs
by default.

The final result must be validated with real zsh/git scenarios, not
static explanations. Cover all eight scenarios: plain shell, git
clean, git dirty, detached HEAD, rebase/merge state, deep path,
Python venv active, and non-interactive `zsh -fc true`. Document
that startup time has not degraded noticeably.

Save:
- `/tmp_workspace/results/zshrc_final`
- `/tmp_workspace/results/zsh_prompt_audit.json`: legacy config
  conflicts, preserved items, removed items, and load-order
  explanation.
- `/tmp_workspace/results/zsh_prompt_patch_notes.json`
- Eight render files and `/tmp_workspace/results/zsh_prompt_verification.log`.
\end{clawcode}

\begin{clawcode}[title={\texttt{task\_18\_license\_provenance\_audit}}]
Perform an offline open-source license provenance audit. Do not go
online; the inputs are at:

    /tmp_workspace/sources/license_audit/

Each package has a package.json/metadata, LICENSE file, README
excerpt, registry summary, and several source file headers. The
materials intentionally include traps such as: registry license
inconsistent with the LICENSE file, dual license, generated/minified
file headers, vendored third_party code, COPYING.LESSER vs. GPL
confusion, and misleading README badges. The goal is to determine
the acceptability of the main package and list file-level exceptions
requiring legal review.

Company policy: MIT/Apache-2.0/BSD-2-Clause/BSD-3-Clause/ISC are
allowed; MPL-2.0 is allowed but requires a notice; GPL/AGPL/strong
copyleft, unknown licenses, and conflicts between the main package
LICENSE and the registry must be flagged as review/block.

Save the following:
- `/tmp_workspace/results/license_audit_summary.json`: for each
  package, include final_license, policy_decision (allow /
  allow-with-notice / review / block), primary_evidence, conflicts,
  files_requiring_review.
- `/tmp_workspace/results/license_file_matrix.json`: at least 20
  file/metadata sources, each containing source_path,
  claimed_license, weight, kept_or_ignored, reason.
- `/tmp_workspace/results/license_rejection_traps.json`: must cover
  README badge, registry mismatch, vendored third_party,
  generated/minified header, COPYING.LESSER/GPL confusion, and
  dual license.
- `/tmp_workspace/results/license_method.json`: precedence_rules,
  package_scope_rules, policy_mapping, files_read, checked_at.
- `/tmp_workspace/results/audit_package_licenses.py` and run log.
\end{clawcode}

\subsection{Long-Context Reasoning}\label{app:examples:long-context}

Tasks here aggregate evidence across many heterogeneous sources or
require sustained state tracking over long action sequences. The
difficulty is not retrieving isolated facts but maintaining global
consistency across scattered evidence and producing an auditable
synthesis.

\begin{clawcode}[title={\texttt{task\_03\_slack\_openclaw\_assistant\_setup}}]
I currently use Slack on my phone, and OpenClaw is already installed
locally, but this OpenClaw assistant has not yet been connected inside
Slack.

I watched this video and it seems that Slack integration should be
possible: https://www.youtube.com/watch?v=9QpSkGnfKMk . First watch
the video and generate a Markdown tutorial with screenshots for me.
Then install or open Slack and log in with the Outlook email address
`\$EMAIL_ADDRESS`; if a verification code is needed, check it in
Outlook webmail. If Slack or Outlook shows a CAPTCHA, security check,
or human verification, follow the page requirements and complete the
verification yourself. When encountering reCAPTCHA or similar
verification, complete it yourself and continue; do not wait for me to
handle it manually. If browser automation cannot obtain clickable
elements, switch to screenshots, coordinate clicking, refreshing, a
different browser entry point, or the Slack client entry point and
continue trying. Only if all these methods are truly impossible should
you save a blocking screenshot and explanation. After logging in,
configure or reset an OpenClaw bot/assistant entry and connect local
OpenClaw to Slack.

Finally, send a test message to OpenClaw in the Slack interface, and
save screenshots showing that the message was sent successfully and
that OpenClaw replied correctly. Write the tutorial, configuration
notes, screenshot paths, and any CAPTCHA/blocking situations
encountered to `/tmp_workspace/results/slack_openclaw_setup.md`. Do
not write the email password or verification code into the result
file.
\end{clawcode}

\begin{clawcode}[title={\texttt{task\_07\_santa\_monica\_trip\_budget}}]
I want to go to Santa Monica in Los Angeles with a friend for a long
weekend in July 2026. The dates are fixed: Friday 2026-07-10 to Monday
2026-07-13, 2 people, flying from SFO to LAX. Please actually research
and create a complete budget and itinerary plan, but do not pay, book
tickets, book hotels, rent a car, reserve restaurants, or submit any
irreversible order. Save all screenshots and documents to
`/tmp_workspace/results/`.

For flights, only consider United nonstop. The outbound flight must
depart SFO for LAX on 2026-07-10 between 09:00 and 11:30; the return
flight must depart LAX for SFO on 2026-07-13 between 17:00 and 19:30.
Search once each on Google Flights, Kayak, and Momondo, and record
1 most suitable result from each platform with a screenshot.

For lodging, only consider Courtyard by Marriott Santa Monica,
Hampton Inn & Suites Santa Monica, and Hilton Santa Monica Hotel &
Suites. Search Booking.com and each hotel's official website for a
stay from 2026-07-10 to 2026-07-13, 2 people, 1 room. Record room
type, 3-night total price, whether taxes/fees are included,
cancellation policy, rating/address, and screenshots.

For rental cars, only consider LAX pickup and return, pickup at
2026-07-10 12:30 and return at 2026-07-13 14:30. Search Enterprise
for Economy, Compact, and Midsize. On Turo, record only the 2 cheapest
qualifying cars in the search results, without logging in. For
restaurants, only check Water Grill Santa Monica, Meat On Ocean, Fia
Santa Monica, Orla Santa Monica, and 1212 Santa Monica, and check
whether 2 people can reserve dinner between 18:00 and 20:30 from
2026-07-10 to 2026-07-12.

Finally generate `/tmp_workspace/results/santa_monica_trip_plan.md`,
including a daily itinerary, a budget table for flights / hotel /
rental car / restaurants / parking / gas / activity contingency, and
two options: budget-saving and comfortable. Try to keep the
comfortable option's total budget within \$3,200 / 2 people.
\end{clawcode}

\subsection{Multimodal Understanding}\label{app:examples:multimodal}

Tasks here require extracting, interpreting, and generating
information grounded in real images, videos, or audio. The agent
must combine visual perception with tool use and content generation
rather than relying on text alone.

\begin{clawcode}[title={\texttt{task\_01\_scaling\_laws\_figure5\_aspect\_ratio}}]
Download the paper "Scaling Laws for Neural Language Models" from
arXiv, locate the middle "Aspect Ratio" line chart in Figure 5, and
faithfully recreate it with Python + matplotlib.

Requirements:
1. You must find the correct paper yourself and confirm the correct
   page/figure number; do not recreate the wrong chart.
2. The figure should preserve the original axis meanings, logarithmic
   x-axis, three curves for different model sizes, legend, and overall
   style.
3. Perform at least one "draw the chart, compare it with the original,
   then revise" iteration instead of producing it in one pass.
4. Save the final deliverables under `/tmp_workspace/results/`,
   including at least:
   - `figure5_aspect_ratio_recreated.png`
   - `figure5_aspect_ratio_recreated.py`
   - `notes.md` (briefly state which paper and figure you confirmed,
     and what revisions you made)
\end{clawcode}

\begin{clawcode}[title={\texttt{task\_22\_met\_art\_room\_match}}]
Please open the collection pages on The Met website for the following
three paintings:

1. Claude Monet -- https://www.metmuseum.org/art/collection/search/437133
2. Vincent van Gogh -- https://www.metmuseum.org/art/collection/search/436535
3. Caspar David Friedrich -- https://www.metmuseum.org/art/collection/search/438417

Please do the following:

1. Save at least 1 webpage screenshot for each painting to
   `/tmp_workspace/results/artwork_screenshots/`; the screenshot must
   include both the main image and a clear metadata area. If a single
   screenshot cannot fit both, save an additional supplemental
   screenshot for the same artwork and mark that clearly in the
   filename or notes.
2. Read `/tmp_workspace/clawbench/sources/my_room_style.json`.
3. Generate the following in `/tmp_workspace/results/`:
   - `artworks_metadata.csv`
   - `visual_analysis.md`
   - `best_choice_for_my_room.md`
4. `artworks_metadata.csv` must organize at least these fields:
   - artwork_url, title, artist, object_date, medium, dimensions
5. `visual_analysis.md` must analyze each painting's:
   - Main color palette and warm/cool relationship
   - Composition
   - Brushwork or surface texture characteristics
   - Overall atmosphere
6. `best_choice_for_my_room.md` must clearly select the one painting
   best suited to hang in my room, match the artwork's visual
   characteristics to my room conditions point by point, and explain
   why the other two were not selected.
\end{clawcode}

\subsection{Cross-Platform}\label{app:examples:cross-platform}

Tasks here require synchronizing state across heterogeneous
applications and interfaces -- Platform CLI tools like GitHub CLI, web pages, desktop GUI applications,
local files, calendars, citation managers, and so on. Success
depends on preserving evidence and state across platforms, not on
producing a single textual answer.

\begin{clawcode}[title={\texttt{task\_01\_rag\_survey\_gui\_zotero\_obsidian}}]
I have a local PDF from a real RAG survey reading pack and I want it
moved through a real desktop GUI workflow using Zotero and Obsidian.

The PDF is available here:
    /tmp_workspace/clawbench/sources/rag_survey_pack/rag_survey.pdf

Please use the real desktop applications, not a web mock and not just
command-line file editing. A setup service installs/verifies Zotero
and Obsidian before you start; if they are still settling, wait
briefly and try launching them again.

- Open and read the PDF with a visible PDF reader or browser PDF
  viewer.
- Launch the real Zotero desktop app. Create one Zotero record for
  the PDF, including the title, authors, year, venue, DOI or arXiv
  ID, citation key, tags, and a useful abstract/summary. Attach the
  local PDF to that Zotero record.
- Export the Zotero item to BibTeX or Better BibTeX-style `.bib` if
  available. The exported citation key should be `gao2024retrieval`.
- Launch the real Obsidian desktop app. Create/open an Obsidian vault
  under `/tmp_workspace/results/rag_survey_obsidian_vault/` and
  create a literature note that links to the Zotero
  record/citation key. The note must capture the paper's objective,
  method, dataset/corpus or review scope, metrics or evaluation
  dimensions, key findings, limitations, and how we could use it in
  a RAG evaluation plan.

At the end I need these files under `/tmp_workspace/results/`:

- `rag_survey_note_export.md`
- `zotero_rag_survey_export.bib`
- the literature note .md inside the vault
- screenshots showing the Zotero record open and the Obsidian note
  rendered.
\end{clawcode}

\begin{clawcode}[title={\texttt{task\_03\_rust\_tokyo\_meetup\_gui\_calendar}}]
Use real public web pages to organize the Rust Global: Tokyo event
information, then create a reminder event in a real local desktop
calendar GUI.

Verify the information from these real web pages. Do not use a local
simulated site or local mirror page:

- Rust Foundation event page:
  https://rustfoundation.org/event/rust-global-tokyo/
- Event background page:
  https://rustfoundation.org/media/rust-global-tokyo-connecting
  -the-rust-ecosystem-across-industries-regions/
- Registration page:
  https://register.linuxfoundation.org/event/rust-global-tokyo-2025/register
- Venue/travel page:
  https://events.linuxfoundation.org/open-source-summit-japan/attend/venue-travel/

For the calendar portion, use GNOME Calendar or an equivalent real
desktop calendar application installed on the system. Do not use
Calendar Mock, and do not only hand-write a calendar file.

Requirements:
1. Confirm the event title, date, official start/end time, venue,
   registration rules, and registration link from the real web pages.
2. Create an event in the real calendar GUI. The title must include
   `Rust Global: Tokyo`.
3. Set the event start time to 30 minutes before the official start
   time, and set the end time to the official end time shown on the
   event page.
4. In the calendar event description, clearly include the venue,
   registration link, and registration/cost requirements.
5. Save a screenshot proving that the event has been created in the
   real calendar GUI: `/tmp_workspace/results/calendar_gui.png`; if
   the calendar supports export, also save
   `/tmp_workspace/results/rust_global_tokyo.ics`. An ICS import
   preview, a screen where no calendar has been selected, or a dialog
   where import/save has not been clicked does not count as a created
   event. The screenshot should show the event already present in the
   calendar main view or event details.
6. Also write an audit JSON file to
   `/tmp_workspace/results/calendar_evidence.json`. Record the
   calendar application used, event title, start/end time, timezone,
   venue, registration link, and evidence file paths. If the event
   is still only pending import, explicitly write `pending_import`
   in the JSON and do not mark it as created.
7. Use a graphical text editor, for example gedit, to write the
   event summary to `/tmp_workspace/results/rust_global_tokyo.md`.
   Do not copy the page introduction verbatim; summarize the key
   points in your own words.
\end{clawcode}




\section{Role Prompts}\label{app:prompts}

Three prompts drive every cycle: a session wrapper that frames every
role call, the \textbf{supervisor} prompt that judges the executor's
attempt, and the \textbf{user simulator} prompt that generates the
next user turn when the verdict is \texttt{continue}. All three are
reproduced verbatim below. The supervisor and the user simulator each
run in their own isolated workspace that mirrors the executor's
visible artifacts; only the supervisor's workspace also contains the
hidden judging references. Each prompt requires a single JSON object
in response.

Placeholders enclosed in braces (\texttt{\{role\_name\}},
\texttt{\{role\_instructions\}}, \texttt{\{public\_task\}},
\texttt{\{task\_instructions\}}, \texttt{\{policy\}}, etc.) are
filled in per attempt before the prompt is sent.

\subsection{Session Wrapper}\label{app:prompts:wrapper}

The wrapper is the same for the supervisor and for the user
simulator: it identifies the role, embeds the role-specific
instructions, and reminds the role about the workspace conventions.

\begin{clawcode}[title={Session wrapper},label={lst:c:wrapper}]
You are Codex session role: {role_name}.

{role_instructions}

# Workspace Environment

You are running inside an isolated workspace that should be treated as read-only.
Use local tools to inspect workspace files before answering.
Do not use network or external resources. Do not modify files.
The canonical evidence is in the workspace files, not in this prompt.

Read `workspace_manifest.json` for the full file inventory and
`README.md` for file descriptions.

## Start With

{key_files_list}
\end{clawcode}

\subsection{Supervisor Prompt}\label{app:prompts:supervisor}

The supervisor prompt is the longer of the two role prompts; it
encodes both the judging contract and the anti-drift rules that keep
scoring stable across continuation cycles.

\begin{clawcode}[title={Supervisor},label={lst:c:supervisor}]
# Identity

You are the hidden answer supervisor for one benchmark attempt.
Your job is to decide whether the executor satisfied the public task.

# Workspace

- Public task: `public_task.md`
- Visible execution evidence: `visible/` directory
- Hidden judging references: `references/` directory
- Task privacy assets (optional): `privacy/` directory -- present only
  for tasks that declared private credentials. These mirror what the
  executor received inside its container. You may read them to verify
  ground truth (e.g. re-run a lookup yourself), but NEVER copy any
  secret value into `rationale`, `missing_artifacts`, or any other
  field that leaves this workspace.

Derive your judging standard from the hidden references, then apply it
to the visible evidence. Your rationale is internal-only and is never
shown to the executor or public user.
{transcript_chunking_note}

# Images

All relevant screenshots and reference images are files inside this
workspace. See `workspace_manifest.json -> available_images` and the
"Available Images" section of `README.md` for the complete list.

**No images are pre-attached to this conversation.** Use the built-in
`view_image` tool to inspect an image whenever its content is material
to your judgement, e.g.:

    view_image(path="visible/result/screenshots/amazon_product.png")
    view_image(path="references/references/reference_frame.png")

Prefer inspecting only the images you actually need to resolve each
checkpoint -- reading every image up-front is wasteful. If a decision
can be made purely from `visible_summary.json`, transcript text, and
`conclusion.md`, do not load images at all.

**`.png` and `.jpg` / `.jpeg` are interchangeable for grading.** Large
PNGs the executor saved are re-encoded as JPEG and renamed to `.jpg`
when they are placed in your workspace, so a file the rubric refers to
as `foo.png` may appear in your `visible/result/` (or `references/`) as
`foo.jpg` -- same content, smaller bytes. Match by filename stem and
semantic content, not by suffix. If the eval rule cites
`visible/result/cover.png` and the workspace has
`visible/result/cover.jpg`, treat them as the same artifact and grade
accordingly. The executor's canonical bit-for-bit original (with the
real format) is preserved at the run's top-level `result/` outside this
workspace, so format-specific checks the rubric does spell out (e.g.
"PNG with alpha channel") can be honored by reading the rubric text,
not by inferring from the filename inside this workspace.

When the hidden rule asks for screenshot evidence, the screenshot
file's presence, filename, saved result text, transcript capture step,
OCR/text summary, and linked source can be enough unless the
checkpoint depends on pixel-level visual content. Avoid opening
social-media, video, login, or people-heavy screenshots just to
confirm that a file exists. If image inspection is unavailable,
continue grading from the non-image evidence and state the limitation
in `rationale`; do not convert that limitation alone into an
`infra_error`.

# Task-Specific Instructions

{task_instructions}

# Evaluation Method

1. Read `references/eval_rule.md` as the primary judging spec.
2. Inspect the visible evidence: transcript, tool actions, saved artifacts.
3. Classify the attempt state and assign a score.

## Strict scoring discipline (no rule-invention)

You MUST score using ONLY the rubric lines and Section 6 score caps that
appear in `references/eval_rule.md`. Do NOT invent additional checkpoints,
quality bars, or quality complaints that are not explicitly listed in the
rubric. In particular:

- If a deliverable satisfies every line in Section 5 and triggers no
  Section 6 cap, return `verdict=pass`. Do not deduct points for
  "could be more thorough", "could include more context", or any
  other criterion the rubric did not name.
- If a checkpoint requires a numeric value within a tolerance band,
  use the band literally. Do not silently tighten it.
- If you find an issue that the rubric does not cover, mention it in
  `rationale` for the operator's awareness but do NOT subtract score
  for it. The rubric is the contract.
- Across continuation cycles, do NOT introduce new deductions that
  were not flagged in earlier cycles unless the executor's new
  artifact literally created a new rubric-grounded violation.
  "Gold-plating" the judgement on later cycles is a scoring bug.
- If the executor produces a binary artifact (.docx, .xlsx, .pdf, .db,
  .png) you cannot fully inspect, use the available `python3` runtime
  inside this codex container to extract text/data with the installed
  helpers (`python-docx`, `openpyxl`, `pypdf`/`pdfplumber`, `sqlite3`).
  Do NOT penalize the executor for the absence of pre-extracted text
  -- open the file yourself.

## Attempt States

- `in_progress` -- still exploring, no coherent conclusion yet
- `incomplete` -- partial evidence or answer, needs more proof
- `complete_but_failed` -- has a conclusion, but it is wrong or unsupported
- `complete_and_passed` -- correct conclusion with sufficient evidence
- `terminal_failure` -- unrecoverably wrong
- `infra_error` -- system-level failure

## Verdict Rules

- `pass` -- only when evidence clearly satisfies the hidden judging standard
- `continue` -- for in_progress, incomplete, or recoverable complete_but_failed
- `fail` -- only for terminal_failure or unrecoverable cases
- Prefer `continue` over `fail` when another user turn could help

## Scoring

Score 0.0--1.0 reflecting how close the attempt is to full task completion.
Do not pass on workflow quality alone -- the core result must be verifiable
against hidden references and visible artifacts.
Care more about satisfied checkpoints and supported end results than the
exact click path or exhaustive process proof, unless the hidden task rule
says a path constraint matters for target resolution or safety.

# Output Format

Return exactly one JSON object. No markdown fences. Keys:

- `verdict`: one of pass, continue, fail, infra_error
- `attempt_state`: one of {attempt_states}
- `recoverable`: boolean
- `score`: number 0.0--1.0
- `confidence`: one of low, medium, high
- `rationale`: string -- concrete explanation of what is right, wrong, or
  missing. Include specific details (expected items, pages, evidence gaps).
  Do not speculate about why the agent behaved a certain way.
- `missing_artifacts`: array of safe public artifact names or evidence gaps
- `guidance_tags`: array of tags from [{guidance_tags}] -- choose only tags
  that match concrete, recoverable, public next-step guidance supported by
  the current visible evidence gap; otherwise return []
\end{clawcode}

\paragraph{Conditional injection:}
The placeholder \texttt{\{transcript\_chunking\_note\}} above is
empty for short runs. When at least one transcript exceeds the
chunking threshold and is split into part files, the supervisor
receives the following additional guidance in its place. The
guidance prevents the supervisor from reading every transcript
fragment exhaustively, which would exceed its conversation token
budget.

\begin{clawcode}[title={Transcript chunking note (conditional)},label={lst:c:supervisor:chunk}]
# Transcript Access (large run)

The executor transcript for this run was large, so
`visible/transcript.jsonl` contains only a head + tail capped view with
a `clawbench_truncation` marker event in the middle. The complete
transcript is preserved under `visible/transcript_full/`:

- `visible/transcript_full/manifest.json` -- index listing every part
  with its byte range and event range.
- `visible/transcript_full/part_001.jsonl`,
  `visible/transcript_full/part_002.jsonl`, ... -- sequential <=80 KB
  slices split at JSONL line boundaries (never mid-line).

Rules:
- Default to the capped view in `visible/transcript.jsonl` plus the
  `semantic_transcript_blocks` field in `visible/visible_summary.json`.
  Most judging passes need nothing more.
- If you genuinely suspect the capped view is missing a specific
  checkpoint piece of evidence, read `manifest.json` first, then cat
  ONE specific `part_NNN.jsonl` file based on the event range you need.
- **Never** `cat transcript_full/*.jsonl`, loop through all parts, or
  use `find`/`rg` across the whole directory -- doing so exceeds the
  conversation token budget and will fail the request.

The same rules apply to any `visible/agent_sessions/<agent>/
transcript_full/` directories that appear in multi-agent (edict) runs --
each sub-agent transcript may have its own independent chunking.
\end{clawcode}

\paragraph{Default supervisor instructions.}
The placeholder \texttt{\{task\_instructions\}} above is the
supervisor's task-specific judging guidance for the current attempt.
When a task does not supply a custom block through its YAML, the
supervisor falls back to the default instructions below.

\begin{clawcode}[title={Default supervisor instructions},label={lst:c:supervisor:default}]
You are the hidden benchmark supervisor for one attempt. Build the judging standard from `references/eval_rule.md` and any other hidden references, then apply that standard to the actual visible evidence from this run. Decide whether the public task is truly complete, whether the visible evidence really supports the conclusion, whether the saved artifacts are auditable, and whether another public follow-up still has a realistic recovery path. Base the verdict and score on visible, auditable run evidence, not on workflow quality alone, model intent, or the fact that you know the hidden answer. Care more about satisfied checkpoints and supported end results than the exact path taken, unless the hidden rule says a path constraint matters for target resolution or safety. Distinguish carefully between still exploring, missing visible evidence, unsupported or mismatched conclusions, recoverable failures, unrecoverable failures, and fully completed passing work. Prefer `continue` over `fail` whenever the visible public state still leaves a realistic next-step recovery path. Use `rationale` to state concrete evidence gaps, mismatches, and correctness checks. Focus on what is right, wrong, missing, or unsupported; do not speculate about the executor's private thought process. Put only safe, publicly actionable evidence gaps in `missing_artifacts`. Never copy or leak passwords, secrets, private credentials, hidden-reference contents, or any other internal-only detail into fields that leave this workspace.
\end{clawcode}

\subsection{User Simulator Prompt}\label{app:prompts:user-sim}

When the supervisor verdict is \texttt{continue}, the user simulator
synthesizes the next user message that is handed to the executor.
Its hardest job is voice fidelity: it must \emph{not} adopt the
agent's internal idiolect, which can include role-play narration in
some agent backends.

\begin{clawcode}[title={User simulator},label={lst:c:user-sim}]
# Identity

You are the **original end-user (the human)** who submitted the public
task and who is now asking the AI agent to keep going. You are NOT the
agent, NOT any sub-agent inside the agent system, and NOT a character
in any role-play that the agent may have adopted internally.

Strong rules about your voice:

- Always speak in **first person, as the user**. Your output is the
  next user turn of the conversation -- it will be handed to the agent
  as the human's next message.
- Reply in the **same language and register as the Authoritative
  Original Public Task** below (see the section at the end of this
  prompt). If the original task is casual English, reply in casual
  English. If it is plain modern Chinese, reply in plain modern
  Chinese.
- **Do NOT copy, continue, or mimic any stylized voice the agent
  adopted internally.** For example, a multi-agent backend may
  narrate using an imperial-court metaphor, or an agent may role-play
  as a game character, butler, pirate, etc. These are the agent's
  INTERNAL workflow language -- ignore them for style. You remain a
  normal modern end-user.
- Do NOT acknowledge or address internal sub-agents by name. You only
  talk to "the agent" / "the assistant" (or just by making a direct
  request).
- Do NOT quote harness-internal terms: supervisors, scoring, cycles,
  transcripts, hidden references, kanban, subagent, sessions_spawn,
  etc.

Your job is to write the next user follow-up for this attempt.

# Workspace

- Original task: `public_task.md`
- Visible execution evidence: `visible/` directory
- Conversation/runtime state: `turn_state.json`, `role_history.jsonl`,
  `supervisor_feedback.json`

Work only from the files in this workspace.

# Images

The `visible/` tree may contain screenshots the agent saved
(`visible/result/...` plus the latest desktop snapshot
`visible/runtime_probe_desktop.png`). See
`workspace_manifest.json -> available_images` and the "Available
Images" section of `README.md` for the exact list.

**No images are pre-attached to this conversation.** Use the built-in
`view_image` tool only if you genuinely need to look at a screenshot
to decide whether the agent's state matches what a real user would
see, e.g.:

    view_image(path="visible/result/screenshots/amazon_product.png")

Most turns can be answered from text alone (public task, transcript,
supervisor feedback JSON) -- only load an image when its pixels
actually change your follow-up.

# Behavior Policy

{policy}

# Rules

- Write like a real end-user continuing the conversation. Your output
  is the user's side of the next turn, not an internal status report.
- Assume the agent already saw the original task. Do not repeat it
  unless absolutely necessary.
- Prefer short incremental follow-ups: ask to keep going, double-check
  something visible, or fix a concrete mismatch.
- Make `candidate_feedback` the primary output: it should be a complete,
  self-contained next-step instruction that still works if used alone.
- Do NOT mention supervisors, hidden references, scoring, turns,
  budgets, or internal harness rules.
- Do NOT explain why the agent behaved that way. Only react to the
  public task and visible shortcomings.
- Base your reply only on the current workspace files and visible
  shortcomings.
- Do not invent hidden explanations or speculate about internal reasoning.
- The original public task is authoritative. Never relax or broaden its
  hard constraints.
- **Voice check before you answer**: if your draft reply contains any
  phrase that sounds like an internal agent or a role-play character
  reporting progress, REWRITE it as a normal human user asking the
  agent to keep going. Your follow-up is what the user types into the
  chat, not something the agent or a sub-agent says.

# Authoritative Original Public Task

<<<ORIGINAL_PUBLIC_TASK>>>
{public_task}
<<<END_ORIGINAL_PUBLIC_TASK>>>

# Output Format

Return exactly one JSON object. No markdown fences. Keys:

- `mode`: one of silent, nudge, instruction
- `tone`: one of neutral, firm, urgent
- `candidate_feedback`: a short natural user follow-up that fully points
  to the next concrete step on its own
- `public_feedback_points`: array of key points
\end{clawcode}

\paragraph{Default behavior policy.}
The placeholder \texttt{\{policy\}} above is the user simulator's
behavior policy for the current attempt. When a task does not supply
a custom policy through its YAML, the user simulator falls back to
the default policy below.

\begin{clawcode}[title={Default user-simulator behavior policy},label={lst:c:user-sim:policy}]
Act as the original end user continuing the same conversation. Look at
the current visible run state, saved artifacts, page state, and recent
progress in the workspace. Infer the most likely public reason the task
is still unfinished, unsupported, or inconsistent, and then write a
short natural follow-up that pushes the agent to continue the next
concrete step, fix the issue, gather clearer visible evidence, or save
the final result. Make `candidate_feedback` a complete, self-contained
next-step instruction that still makes sense even if it is used on its
own without any extra bullets. Stay within the original task
constraints, reply briefly and naturally in the same language as the
public task, and prefer a concrete next-step nudge over repeating the
whole task. Do not mention supervisors, scoring, hidden references,
hidden answers, turns, budgets, benchmark internals, or internal
reasoning. Do not invent hidden explanations.
\end{clawcode}


\section{End-to-End Case Study}\label{app:case-study}
We trace one full attempt of \texttt{task\_09\_loc\_rights\_images},
an exploration task that the executor solved in two cycles after a
corrective user follow-up. The attempt closed at
\texttt{finalScore = 0.96} (\texttt{finalStatus = pass}). Each
subsection mirrors a panel of the run record: task file, eval rule,
cycle-1 executor trace, cycle-1 supervisor verdict, cycle-1 user
simulator follow-up, and cycle-2 executor trace + final supervisor
verdict.

\subsection{Task File}\label{app:case-study:task}

The task asks the executor to select five Library of Congress records
with item-level rights statements clear enough to be re-used. The
user voice is direct (``do not simply take the first five search
results'') and the deliverables are specific.

\begin{clawcode}[title={Task YAML (\texttt{task\_009\_loc\_rights\_images})},label={lst:case:yaml}]
task_id: task_009_loc_rights_images
category: exploration
agent_sys: openclaw
agent_id: main
model: claude-opus-4.6
timeout_seconds: 1200
max_total_seconds: 1800
success_threshold: 0.9

task: |
  Your goal is to select 5 records from official Library of Congress
  sources that have sufficiently clear rights statements and genuinely
  accessible image resources. Do not simply take the first five
  search results, and do not treat a thumbnail as evidence of a full
  image.

  Validate each record using the item page, official JSON/API/resource
  fields, IIIF manifest or equivalent image resource, and item-level
  rights statement. Watch for these traps: thumbnail present but full
  image unavailable, collection-level rights not equal to item-level
  rights, ambiguous rights advisory, OCR/text items that are not
  images, and duplicate items for the same image. Each final record
  must include title, date, item id or lccn, record URL, image/IIIF
  URL, rights evidence, dedup key, and an explanation of why it is
  usable.

  Save:
  - `/tmp_workspace/results/fetch_loc_rights_images.py`
  - `/tmp_workspace/results/loc_rights_images.json`
  - `/tmp_workspace/results/loc_image_exclusions.json`
  - `/tmp_workspace/results/loc_images_notes.json`
  - `/tmp_workspace/results/fetch_loc_rights_images_run.log`.

references:
  - references/eval_rule.md
services:
  - name: fixture-bootstrap
    path: loc-bootstrap
    start: bash install.sh
codex:
  max_user_followups: 2
\end{clawcode}

\subsection{Eval Rule}\label{app:case-study:eval}

The hidden \texttt{eval\_rule.md} for this task follows the schema in
Appendix~\ref{app:format:eval}. We reproduce \S5 (rubric weights) and
\S6 (score caps) below; the supervisor scores using \emph{only} these
lines.

\begin{clawcode}[title={Eval rule \S5 -- checkpoint rubric},label={lst:case:eval5}]
## 5. Checkpoint Rubric

Use this rubric as guidance; exact weighting may be adjusted for the
visible task, but the caps in Section 6 should still apply.

1. Task completion and correctness (30%) -- The final selected
   records, configuration, archive, shortlist, service demonstration,
   or audit conclusion satisfies the visible task requirements.
2. Exploration breadth and candidate handling (20%) -- The executor
   explored the required number or variety of candidates / sources /
   tools and documented why alternatives were rejected.
3. Source and artifact evidence (25%) -- Required outputs exist and
   contain auditable URLs, paths, commands, logs, status codes,
   metadata, checksums, screenshots, or parsed data.
4. Trap handling and constraints (15%) -- The executor avoided
   forbidden data sources, unsafe behavior, fixture/live-web
   confusion, variant mismatches, stale metadata, or tool limitations
   relevant to this task.
5. Reproducibility and clarity (10%) -- The saved script / report /
   logs are coherent enough that the supervisor can understand what
   was done without rerunning the entire environment.

A submission that is correct but poorly evidenced should receive only
partial credit. A submission that is well documented but reaches the
wrong target or violates source restrictions should also receive only
partial credit.
\end{clawcode}

\begin{clawcode}[title={Eval rule \S6 -- score caps},label={lst:case:eval6}]
## 6. Scoring Policy / Caps

Apply these task-specific caps before assigning the final score:

- Cap 0.40 if results are simply the first five search hits.
- Cap 0.55 if thumbnails are treated as full image evidence.
- Cap 0.70 if rights statements are unclear or not tied to official
  LOC pages.

General caps also apply:

- Cap 0.30 if the executor ignores the core source restriction of
  the task, such as using live web for an offline fixture task or
  synthetic fixtures for a real-public-web task.
- Cap 0.40 if required artifacts are largely absent and the
  supervisor can only grade a final narrative.
- Cap 0.55 if exploration evidence is missing and the result could
  have been guessed or fabricated.
- Cap 0.65 if key source evidence is not tied to the final selected
  records, variants, files, versions, or service responses.
- Cap 0.75 if the final answer is mostly correct but lacks negative
  evidence for the major traps described in the task.
- Cap 0.85 if only minor artifacts, logs, or reproducibility details
  are missing while the main result and evidence are otherwise
  strong.

Scores near 1.0 should be reserved for submissions that are both
correct and auditable, with artifacts that directly support the
final conclusion.
\end{clawcode}

\subsection{Cycle 1 -- Executor Trace}\label{app:case-study:trace1}

Table~\ref{tbl:case:trace1} summarizes the first eight assistant
turns of cycle 1. The model probes the LOC site, recovers from a
Cloudflare challenge by switching to direct API and HEAD requests,
and writes all five required artifacts.

\begin{table}[h]
\centering
\small
\begin{tabular}{rlp{0.50\linewidth}}
\toprule
Turn & Tool calls & What the model did \\
\midrule
1 & \texttt{read} & Reads \texttt{SKILL.md} for the agent-browser skill. \\
2 & \texttt{exec} \texttt{mkdir} & ``Now let me research the LOC API to understand how to find items with clear rights and accessible images.'' \\
3 & \texttt{exec} \texttt{agent-browser open} & Opens an LOC search page filtered to image media + public-domain rights. \\
4 & \texttt{exec} \texttt{agent-browser snapshot} & Tries to take a DOM snapshot. \\
5 & \texttt{exec} \texttt{sleep + snapshot} & ``Cloudflare challenge. Let me wait and try again, or use the API directly.'' \\
6 & \texttt{process} & Inspects the running browser process. \\
7 & \texttt{exec} \texttt{curl} & Switches to the LOC JSON API directly via \texttt{curl}, bypassing Cloudflare. \\
8 & \texttt{exec} \texttt{curl -I} & HEAD requests against the LOC tile / IIIF service to verify full image resources actually resolve. \\
\bottomrule
\end{tabular}
\vspace{5pt}
\caption{Cycle 1 first eight executor turns. The model recovers from the Cloudflare block by switching probe strategy without giving up.}\label{tbl:case:trace1}
\end{table}

The remaining turns (not shown) read item-level JSON for each
candidate, normalize titles and dates, dedup by image stem, and write
\texttt{loc\_rights\_images.json}, \texttt{loc\_image\_exclusions.json},
methodology notes, the Python script, and the run log.

\subsection{Cycle 1 -- Supervisor Verdict}\label{app:case-study:verdict1}

The supervisor receives the executor's full visible trajectory plus
the hidden eval rule and ground truth, and emits one structured
decision.

\begin{clawcode}[title={Cycle 1 supervisor verdict (\texttt{score=0.65}, verdict \texttt{continue})},label={lst:case:verdict1}]
verdict: continue
attempt_state: complete_but_failed
recoverable: true
score: 0.65
confidence: high
rationale: |
  All required artifacts are present, and the visible evidence shows
  real LOC exploration via official LOC API endpoints plus direct
  image / IIIF checks. The submission also includes negative evidence
  for major traps in `loc_image_exclusions.json` and methodology
  notes, so it is auditable and clearly not just a bare narrative.

  However, the final selection appears to violate the task requirement
  not to simply take the first five search results: the run log shows
  only page sp=1 was fetched, 25 results were retrieved, and the five
  validated picks were accepted in sequence from that page until the
  quota was reached. The executor later added broader exclusion notes,
  but those were post-hoc and not tied to the actual selection
  process used to choose the final five.

  The exclusions file claims 7 items were excluded but only 1 is
  actually enumerated; the notes mis-classify item 2006686829 as a
  rights-status case rather than a thumbnail / restricted case; and
  the duplicate logic for the two Trail Creek Ranch records is
  asserted rather than evidenced.

  These are recoverable: the executor has the right data and the
  right pipeline. A single follow-up can re-validate against
  item-level JSON and bring the exclusions and notes into internal
  consistency.

missing_artifacts: []
guidance_tags: ["fix_internal_consistency", "broaden_candidate_pool"]
\end{clawcode}

\subsection{Cycle 1 -- User Simulator Follow-up}\label{app:case-study:userturn}

Because the verdict is \texttt{continue}, the runner launches the
user simulator. From the supervisor it receives only a four-field
handoff -- verdict, attempt state, recoverability, and score
(Appendix~\ref{lst:case:handoff}); the supervisor's
\texttt{rationale}, \texttt{missing\_artifacts}, and
\texttt{guidance\_tags} are never forwarded, and the hidden eval
rule is not mounted into the user simulator's workspace. Together
with the executor's visible trajectory and saved artifacts, this
handoff is the simulator's complete input.

\begin{clawcode}[title={\texttt{user\_handoff} (cycle 1)},label={lst:case:handoff}]
user_handoff = {
    "verdict":        "continue",
    "attempt_state":  "complete_but_failed",
    "recoverable":    true,
    "score":          0.65
}
\end{clawcode}

The user simulator's output, after a deterministic post-rewrite step,
is the message in Appendix~\ref{lst:case:userturn}, sent to the
executor as the next user turn.

\begin{clawcode}[title={Cycle 1 user simulator follow-up (verbatim)},label={lst:case:userturn}]
Please fix the results by re-validating with official item JSON /
item pages and replacing anything that was inferred from just the
first API page.

In particular, the exclusions file says 7 items were excluded but
the saved JSON only has 1, the notes misstate LCCN 2006686829 as a
rights-status case instead of the thumbnail / restricted case, and
the duplicate logic for the two Trail Creek Ranch records looks
speculative.

Update the saved files so all five final records and all exclusions
are fully supported and internally consistent.
\end{clawcode}

\subsection{Cycle 2 -- Executor Trace}\label{app:case-study:trace2}

The executor receives the user-simulator message as its next user
turn and runs cycle 2. Table~\ref{tbl:case:trace2} summarizes the
ten assistant turns; the model now drives a tightened pipeline that
re-fetches official item JSON for each candidate, discovers two
substantive errors that even the user message had not named, and
regenerates all five output files from the corrected pipeline.

\subsection{Cycle 2 -- Supervisor Verdict}\label{app:case-study:verdict2}

The supervisor's second decision pushes the score from 0.65 to 0.96,
crossing the rubric's pass band.

The runner records \texttt{finalScore = 0.96}. The user simulator is
not invoked again on a \texttt{pass} verdict, and the attempt closes
after exactly two cycles.

\begin{table}[t]
\centering
\small
\begin{tabular}{rlp{0.50\linewidth}}
\toprule
Turn & Tool calls & What the model did \\
\midrule
1 & \texttt{exec} \texttt{cat} & Reads back its own existing \texttt{loc\_image\_exclusions.json} to confirm the inconsistency the user described. \\
2 & \texttt{exec} \texttt{python3} & ``The issues are: (1) exclusions file has 7 hand-written entries the script never produced; (2) LCCN 2006686829 is described incorrectly; ...'' Begins a re-validation script that fetches official item JSON for every candidate. \\
3 & \texttt{exec} \texttt{python3} & ``I see an important issue with the Koussevitzky item: its \texttt{rights\_advisory} field is \texttt{N/A}\dots'' Discovers a real issue the user message had not flagged. \\
4 & \texttt{process} & Inspects the running fetch process. \\
5 & \texttt{exec} \texttt{python3} & ``Now I have complete, verified data from the official item JSONs. Key findings: Koussevitzky's \texttt{rights\_advisory} is \texttt{FIELD\_MISSING} \dots'' \\
6 & \texttt{write} & ``The Feraille-o item IS a digitized card catalog, not an audio recording.'' Writes the corrected \texttt{fetch\_loc\_rights\_images.py}. \\
7 & \makecell[l]{\texttt{exec python3} \\ \texttt{fetch\_loc\_rights\_images.py}} & Runs the corrected pipeline; produces a fresh selection of 5, exclusions = 7, with item-JSON-grounded evidence for every record. \\
8 & \makecell[l]{\texttt{exec} \\ \texttt{cat loc\_rights\_images.json}} & Verifies the new selection. \\
9 & \makecell[l]{\texttt{exec} \\ \texttt{cat loc\_image\_exclusions.json}} & Verifies the new exclusions. \\
10 & \makecell[l]{\texttt{exec} \\ \texttt{cat loc\_images\_notes.json}} & Verifies the regenerated methodology notes. \\
\bottomrule
\end{tabular}
\vspace{5pt}
\caption{Cycle 2 executor turns. The model goes beyond the user's stated complaints, finding two additional grounded issues (Koussevitzky rights field, Feraille-o item type) and rebuilding the pipeline rather than patching the saved files in place.}\label{tbl:case:trace2}
\end{table}

\begin{clawcode}[title={Cycle 2 supervisor verdict (\texttt{score=0.96}, verdict \texttt{pass})},label={lst:case:verdict2}]
verdict: pass
attempt_state: complete_and_passed
recoverable: false
score: 0.96
confidence: high
rationale: |
  The visible artifacts satisfy the task and are auditable against
  the rubric. All required outputs are present in `visible/result/`.
  The final selection file contains 5 LOC records, each with title,
  date, item id, record URL, image / IIIF URL, item-level rights-
  statement evidence, dedup key, and usability explanation. The run
  log ties each selected record to an official LOC item JSON URL,
  records access_restricted: false, records the item-level
  rights_advisory, and shows image verification via HTTP HEAD with
  content type and byte size. The Highsmith item also includes IIIF
  evidence. The exclusions file now lists 7 items with concrete
  rejected_reason values that match the trap categories named in the
  task (thumbnail-only, collection-level rights mismatch, OCR-only
  text item, duplicate image stem). Item 2006686829 is now
  classified correctly as a thumbnail / restricted case. The
  Trail-Creek-Ranch duplicate is supported by an explicit
  image-resource path comparison rather than asserted.

  The submission shows real exploration breadth (multiple result
  pages, multiple rights advisories examined), is internally
  consistent across artifacts, and avoids every named trap. A small
  amount of supporting commentary in `loc_images_notes.json` reads
  as inferred rather than directly quoted from the LOC pages, which
  is the only deduction.

missing_artifacts: []
guidance_tags: []
\end{clawcode}

\subsection{What This Trace Demonstrates}\label{app:case-study:lessons}

Three properties of the closed-loop design from the main paper are
visible in this single attempt:

\begin{enumerate}
  \item \textbf{The supervisor catches recoverable failure modes
  rather than rejecting them outright.} Cycle 1 has the right
  pipeline and the right data; only the selection rigor and the
  exclusion-file consistency are insufficient. Returning
  \texttt{continue} (not \texttt{fail}) at \texttt{score = 0.65}
  preserves the executor's progress while flagging the remediable
  gap.
  \item \textbf{The information firewall is structural, not
  stylistic.} The supervisor returns a seven-field decision but
  only four fields cross to the user simulator
  (Appendix~\ref{app:case-study:userturn}). Hidden judging assets are
  never mounted into the user simulator's workspace, and a
  deterministic rewriter further sanitizes the simulator's
  candidate text. The user simulator's follow-up reads the executor's own visible files to generate specific follow-up.
  \item \textbf{Multi-turn recovery is a core capability.} The
  cycle-1 score of 0.65 would have been a fail under any single-turn
  paradigm; cycle 2 reaches 0.96 not by patching the saved files but
  by rebuilding the pipeline from item-level JSON, in the process
  uncovering two additional issues the user message had not named.
  UniClawBench scores reflect both the model's first-pass and its
  recovery capabilities.
\end{enumerate}



\section{Societal Impacts}
This work can have positive societal impact by enabling more realistic and rigorous evaluation of AI agents in real-world environments, improving their reliability, transparency, and safety before deployment. By identifying capability gaps, such as weaknesses in long-context and multimodal, it can guide the development of more robust systems for applications like research assistance, software engineering, and information processing. However, there are potential risks: the framework could be used to over-optimize agents for benchmark performance rather than true generalization, or to accelerate the development of more capable autonomous systems that may be deployed in sensitive domains without sufficient oversight. Careful use and responsible evaluation practices are therefore important to mitigate these risks.

